\section{Bảng 1 nói ba chuyện, và người ta chỉ trích một}
\label{sec:ch08-bang-mot}

\index{bảng 1!ba cột đo ba thứ khác nhau}%
Mục~\ref{sec:ch07-tu-nhin} chép bảng 1 ra và dùng cột giữa. Mục này đọc cả ba
cột, vì chúng đo ba đại lượng khác nhau và chỉ một trong ba là chuyện của phần
cứng.

Chú thích bảng nói rõ ký hiệu: \(n\) là độ dài chuỗi, \(d\) là số chiều biểu
diễn, \(k\) là kernel của \tn{tích chập}{convolution}, \(r\) là bán kính của
self-attention bị hạn
chế. Mục 4 của bài đặt tên cho ba tiêu chí trước khi lập bảng:

\begin{quote}
\enquote{One is the total computational complexity per layer. Another is the
amount of computation that can be parallelized, as measured by the minimum
number of sequential operations required.}
\end{quote}

Rồi tiêu chí thứ ba, ở đoạn sau: \enquote{The third is the path length between
long-range dependencies in the network.}

Ba tiêu chí ấy không quy về nhau được, và đó là chỗ đáng dừng. Cột
\code{Sequential Operations} là cột chương trước nói về, và nó đi từ \(O(n)\)
xuống \(O(1)\): một đại lượng về \emph{thứ tự phụ thuộc}, tức về việc cái gì
buộc phải chờ cái gì. Cột \code{Maximum Path Length} là đại lượng về gradient,
và chương~\ref{ch:mot-vector-cho-moi-tu} đã dẫn xuất một lần rồi. Chỉ cột
\code{Complexity per Layer} là về khối lượng tính toán, và nó là cột duy nhất
trong ba cột mà một cái đồng hồ có thể phản bác.

\subsection*{Câu bào chữa cho \texorpdfstring{\(n^2\)}{n2}}

\index{bảng 1!câu điều kiện n nhỏ hơn d}%
Bài không giấu ô \(O(n^2 \cdot d)\). Bài đặt cạnh nó một câu điều kiện, và cả
chương này chạy trên câu ấy:

\begin{quote}
\enquote{In terms of computational complexity, self-attention layers are faster
than recurrent layers when the sequence length \(n\) is smaller than the
representation dimensionality \(d\), which is most often the case with sentence
representations used by state-of-the-art models in machine translations, such as
word-piece and byte-pair representations.}
\end{quote}

Câu ấy có hai vế và chúng chịu trách nhiệm khác nhau. Vế đầu là toán:
\(n^2 d < n d^2\) khi và chỉ khi \(n < d\). Vế sau là một phát biểu về thực tế
năm 2017, rằng câu dịch máy thường ngắn hơn bề rộng mô hình, và ở
\(\dmodel = 512\) với câu vài chục token thì đúng là như vậy.

Nhưng vế đầu bỏ hết hằng số, và \(O\) lớn được phép bỏ. Một tầng con
self-attention có bốn phép chiếu trước khi chấm điểm; một tầng Long Short-Term
Memory, viết tắt là LSTM, có bốn cổng.
Bốn với bốn ấy không triệt tiêu nhau, và mục~\ref{sec:ch08-dem} là chỗ viết
chúng ra.

\index{bảng 1!chữ recurrent chỉ hai thứ}%
Chỗ mơ hồ thứ hai nằm ở chính chữ \code{Recurrent}. Bảng viết một hàng cho cả
họ. Trong họ ấy, một tầng hồi quy trơn làm hai phép nhân ma trận mỗi bước, còn
một tầng LSTM làm tám, vì bốn cổng và mỗi cổng hai phép. Cùng một ô \(O(n
\cdot d^2)\), và hệ số 4 nằm giữa hai thành viên. Sách này dùng LSTM từ
chương~\ref{ch:encoder-decoder}, nên khi chương này so, nó phải nói đang so với
cái nào.

\subsection*{Một câu của bài mà tôi chưa thấy ai trích}

\index{bảng 1!tích chập tách được}%
Cùng trang, hai đoạn dưới bảng, bài tự xếp kiến trúc của mình ngang một phép
tích chập:

\begin{quote}
\enquote{Even with \(k = n\), however, the complexity of a separable convolution
is equal to the combination of a self-attention layer and a point-wise
feed-forward layer, the approach we take in our model.}
\end{quote}

Đọc chậm câu ấy thì nó nói: về \tn{độ phức tạp}{complexity}, cái bài dựng không
rẻ hơn một tích
chập tách được với kernel phủ cả câu. Bài in nó ra và không khai thác, vì luận
điểm của bài không phải chuyện rẻ. Luận điểm là cột giữa, tức chuyện không phải
chờ.

Tôi để câu ấy ở đây vì nó chỉnh lại kỳ vọng cho cả chương. Cái Transformer mua
được không phải ít phép tính hơn. Nó là cùng chừng ấy phép tính, xếp lại thành
hình dạng chạy được cùng lúc. Mục~\ref{sec:ch08-dong-ho} là chỗ hỏi cái hình
dạng ấy trên một CPU đáng bao nhiêu, và câu trả lời phụ thuộc vào một thứ không
có trong bảng 1 lẫn trong ba tiêu chí của mục 4.
